\Needspace{10\baselineskip}
\begin{center}\emph{Die Reste $+5$ und $-5$.}\end{center}

\Article{121}
Zahlenbeispiele legen nahe, dass $+5$ ein Nichtrest jeder ungeraden Zahl der Form $5n+2$ oder $5n+3$ ist, \emph{das heißt} jeder ungeraden Zahl, die ein Nichtrest von $5$ ist. Dass diese Regel keine Ausnahme zulässt, beweisen wir so. Die kleinste Ausnahme, falls es eine gibt, sei $=t$. Sie ist also ein Nichtrest von $5$, während $5$ ein Rest von $t$ ist. Sei $aa=5+tu$, wobei $a$ gerade und kleiner als $t$ ist. Dann ist $u$ ungerade und kleiner als $t$, und $+5$ ist ein Rest von $u$. Ist nun $a$ nicht durch $5$ teilbar, so ist es auch $u$ nicht. Aber $tu$ ist offensichtlich ein Rest von $5$, und da $t$ ein Nichtrest von $5$ ist, ist auch $u$ ein Nichtrest. \emph{Das heißt,} es gibt einen ungeraden Nichtrest von $5$, für den $+5$ ein Rest ist, unterhalb von $t$, entgegen der Voraussetzung. Ist dagegen $a$ durch $5$ \SourceLeaf{96}teilbar, so setze man $a=5b$ und $u=5v$. Daraus folgt $tv\equiv -1\equiv 4\,(\operatorname{mod.}5)$, \emph{das heißt,} $tv$ ist ein Rest von 5. Der übrige Beweis verläuft wie im ersten Fall.

\Article{122}
Für alle Primzahlen, die Nichtreste von 5 sind und zugleich die Form $4n+1$ haben, \emph{das heißt} für alle Primzahlen der Form $20n+13$ oder $20n+17$, sind also sowohl $+5$ als auch $-5$ Nichtreste. Für alle Primzahlen der Form $20n+3$ oder $20n+7$ ist dagegen $+5$ ein Nichtrest, $-5$ ein Rest.

Ganz ähnlich lässt sich zeigen, dass $-5$ ein Nichtrest aller Primzahlen der Formen $20n+11$, $20n+13$, $20n+17$, $20n+19$ ist. Daraus folgt, dass $+5$ ein Rest aller Primzahlen der Form $20n+11$ oder $20n+19$ und ein Nichtrest aller Primzahlen der Form $20n+13$ oder $20n+17$ ist. Da jede Primzahl außer 2 und 5 (für die $\pm 5$ ein Rest ist) eine der Formen $20n+1,3,7,9,11,13,17,19$ hat, können wir nun alle Fälle entscheiden, außer denen der Form $20n+1$ oder $20n+9$.

\Article{123}
Zahlenbeispiele legen nahe, dass $+5$ und $-5$ Reste aller Primzahlen der Form $20n+1$ oder $20n+9$ sind. Gilt dies allgemein, so haben wir das elegante Gesetz, dass \emph{$+5$ ein Rest aller Primzahlen ist, die Reste von 5 sind} (denn diese haben eine der Formen $5n+1$ oder $5n+4$, also eine der Formen $20n+1,9,11,19$; für die dritte und vierte dieser Formen ist die Aussage schon bewiesen), \emph{dagegen ein Nichtrest aller ungeraden Zahlen, die Nichtreste von 5 sind}, wie wir bereits oben gezeigt haben. Dieser Satz genügt offensichtlich, um zu entscheiden, ob $+5$ (und damit auch $-5$, als Produkt von $+5$ und $-1$ betrachtet) ein Rest oder Nichtrest einer beliebigen gegebenen Zahl ist. Schließlich beachte man die Analogie zu dem Satz über den Rest $-3$ in Artikel 120.

\Needspace{12\baselineskip}
Diese Beobachtung zu beweisen ist jedoch nicht so leicht. Für eine Primzahl der Form $20n+1$ oder allgemeiner der Form $5n+1$ lässt sich ähnlich wie in den Artikeln 114 und 119 vorgehen. Für den Modul $5n+1$ sei $a$ eine beliebige zum Exponenten 5 gehörige Zahl; solche Zahlen gibt es nach dem vorigen Abschnitt. \SourceLeaf{97}Dann gilt $a^5\equiv1$ oder $(a-1)(a^4+a^3+a^2+a+1)\equiv0\ (\mathrm{mod.}\,5n+1)$. Da aber $a\equiv1$ und damit auch $a-1\equiv0$ unmöglich ist, muss $a^4+a^3+a^2+a+1$ $\equiv0$ sein. Somit ist auch $4(a^4+a^3+a^2+a+1)=(2aa+a+2)^2-5a^2$ $\equiv0$, \textit{das heißt,} $5a^2$ ist ein Rest von $5n+1$, also auch $5$, denn $a^2$ ist ein nicht durch $5n+1$ teilbarer Rest (denn $a$ ist nicht durch $5n+1$ teilbar, da $a^5\equiv1)$. \mbox{\textit{Q. E. D.}}

Der Fall einer Primzahl der Form $5n+4$ erfordert dagegen feinere Methoden. Da wir die Sätze, mit denen er sich entscheiden lässt, später allgemeiner behandeln werden, gehen wir hier nur kurz auf sie ein.

\Needspace{8\baselineskip}
I. Ist $p$ eine Primzahl und $b$ ein gegebener quadratischer Nichtrest von $p$, so ist der Wert des Ausdrucks

\[
(A)\ldots\quad\frac{(x+\surd b)^{p+1}-(x-\surd b)^{p+1}}{\surd b}
\]

(aus dessen Entwicklung die Irrationalität offensichtlich verschwindet) stets durch $p$ teilbar, welche Zahl man auch für $x$ einsetzt. Betrachtet man die Koeffizienten in der Entwicklung von $A$, so sieht man, dass alle Glieder vom zweiten bis zum vorletzten einschließlich durch $p$ teilbar sind. Also gilt $A\equiv2(p+1)(x^p+xb^{\frac{p-1}{2}})$ $(\mathrm{mod.}\,p)$. Da aber $b$ ein Nichtrest von $p$ ist, gilt $b^{\frac{p-1}{2}}\equiv-1\ (\mathrm{mod.}\,p)$ (Artikel $106$), während $x^p$ stets $\equiv x$ ist (voriger Abschnitt). Daraus folgt $A\equiv0$. \mbox{\textit{Q. E. D.}}

\Needspace{8\baselineskip}
II. In der Kongruenz $A\equiv0\ (\mathrm{mod.}\,p)$ hat die Unbestimmte $x$ den Grad $p$, und alle Zahlen $0,1,2....p-1$ sind ihre Wurzeln. Sei nun $e$ ein Teiler von $p+1$. Dann ist der Ausdruck

\[
\frac{(x+\surd b)^e-(x-\surd b)^e}{\surd b}
\]

(den wir mit $B$ bezeichnen) nach der Entwicklung frei von Irrationalität; die Unbestimmte $x$ hat darin den Grad $e-1$. Aus den Grundlagen der Analysis ist bekannt, dass $A$ durch $B$ identisch teilbar ist. Ich behaupte nun, dass es $e-1$ Werte von $x$ gibt, die, in $B$ eingesetzt, $B$ durch $p$ teilbar machen. Setzen wir nämlich $A\equiv BC$, so hat $x$ in $C$ den Grad $p-e+1$. Die Kongruenz $C\equiv0\ (\mathrm{mod.}\,p)$ hat daher höchstens $p-e+1$ Wurzeln. Somit sind alle übrigen Zahlen unter $0,1,2,3....p-1$, deren Anzahl $=e-1$ ist, Wurzeln der Kongruenz $B\equiv0$.

\Needspace{8\baselineskip}
III. Nun habe $p$ die Form $5n+4$, es sei $e=5$, $b$ ein Nichtrest von $p$, und die Zahl $a$ sei so bestimmt, dass

\SourceLeaf{98}\[
\frac{(a+\surd b)^5-(a-\surd b)^5}{\surd b}
\]

durch $p$ teilbar ist. Dieser Ausdruck wird aber zu

\[
=10a^4+20aab+2bb=2((b+5aa)^2-20a^4)
\]

Also ist auch $(b+5aa)^2-20a^4$ durch $p$ teilbar, \textit{das heißt,} $20a^4$ ist ein Rest von $p$. Da aber $4a^4$ ein nicht durch $p$ teilbarer Rest ist (denn $a$ kann offensichtlich nicht durch $p$ teilbar sein), ist auch $5$ ein Rest von $p$. \mbox{\textit{Q. E. D.}}

Damit gilt der am Anfang dieses Artikels angegebene Satz allgemein. ---

Wir bemerken noch, dass die Beweise für beide Fälle von dem bedeutenden Lagrange stammen, \textit{M\'em. de l'Ac. de Berlin}, $1775$, \textit{S.} $352$ ff.

\Needspace{10\baselineskip}
\begin{center}\textit{Über $\pm7$.}\end{center}

\Article{124}
Durch eine ähnliche Methode beweist man, dass

$-7$ \textit{ein Nichtrest jeder Zahl ist, die ein Nichtrest von $7$ ist.}

Zahlenbeispiele legen dagegen nahe, dass

$-7$ \textit{ein Rest jeder Primzahl ist, die ein Rest von $7$ ist.}

Doch dies hat bisher niemand streng bewiesen. Für die Reste von $7$, die die Form $4n-1$ haben, ist der Beweis leicht. Mit einer aus dem Vorigen schon vertrauten Methode lässt sich zeigen, dass $+7$ stets ein Nichtrest solcher Primzahlen ist, also $-7$ ein Rest. Damit ist aber wenig gewonnen, denn die übrigen Fälle lassen sich so nicht behandeln. Einen weiteren Fall können wir ähnlich wie in den Artikeln $119$, $123$ erledigen. Ist nämlich $p$ eine Primzahl der Form $7n+1$ und $a$ für den Modul $p$ eine zum Exponenten $7$ gehörige Zahl, so ist

\[
\frac{4(a^7-1)}{a-1}=(2a^3+a^2-a-2)^2+7(a^2+a)^2
\]

offensichtlich durch $p$ teilbar, also $-7(a^2+a)^2$ ein Rest von $p$. Aber $(a^2+a)^2$ ist als Quadrat ein Rest von $p$ und zudem nicht durch $p$ teilbar. Denn da $a$ nach Voraussetzung zum Exponenten $7$ gehört, kann es weder $\equiv0$ noch $\equiv-1\ (\mathrm{mod.}\,p)$ sein. \textit{Das heißt,} weder $a$ noch $a+1$ ist durch $p$ teilbar und somit auch nicht das Quadrat $(a+1)^2a^2$. Also ist auch $7$ offensichtlich ein Rest von $p$. \SourceLeaf{99}\mbox{\textit{Q. E. D.}} --- Die Primzahlen der Form $7n+2$ oder $7n+4$ entziehen sich jedoch allen bisher angegebenen Methoden. Auch diesen Beweis hat der bedeutende Lagrange zuerst an der angegebenen Stelle gefunden. --- In Abschnitt VII werden wir allgemein zeigen, dass sich der Ausdruck $\frac{4(x^p-1)}{x-1}$ immer auf die Form $X^2\mp pY^2$ bringen lässt (wobei das obere Zeichen gilt, wenn $p$ eine Primzahl der Form $4n+1$ ist, das untere bei der Form $4n+3$); dabei bezeichnen $X,Y$ rationale Funktionen von $x$, die keine Brüche enthalten. Über den Fall $p=7$ hinaus hat Lagrange diese Zerlegung nicht ausgeführt; siehe die angegebene Stelle, S. $352$.

\Needspace{24\baselineskip}
\begin{center}\textit{Vorbereitung auf die allgemeine Untersuchung.}\end{center}

\Article{125}
Da die vorigen Methoden nicht ausreichen, um allgemeine Beweise zu führen, ist es Zeit, eine andere ohne diesen Mangel anzugeben. Wir beginnen mit einem Satz, dessen Beweis unseren Bemühungen ziemlich lange widerstanden hat, obwohl er auf den ersten Blick so offensichtlich erscheint, dass einige nicht einmal die Notwendigkeit eines Beweises erkannt haben. Er lautet: \textit{Jede Zahl außer den positiv genommenen Quadraten ist ein Nichtrest gewisser Primzahlen.} Da wir diesen Satz nur als Hilfsmittel für andere Beweise verwenden, behandeln wir hier nur die dafür nötigen Fälle. Für die übrigen Fälle wird sich dasselbe später von selbst ergeben. Wir zeigen also, dass \textit{jede Primzahl der Form $4n+1$, ob positiv oder negativ genommen}\HistoricalNote{*)}{Dass $+1$ ausgeschlossen werden muss, ist für sich klar.}, \textit{ein Nichtrest gewisser Primzahlen ist}, und zwar, falls $>5$, solcher Primzahlen, die kleiner als sie selbst sind.

\textit{Erstens:} Sei die Primzahl $p$ der Form $4n+1$ ($>17$; aber $-13\,N3$, $-17\,N5$) \textit{negativ} zu nehmen. Sei $2a$ die gerade Zahl unmittelbar oberhalb von $\surd p$. Dann ist $4aa$ offensichtlich stets $<2p$, also $4aa-p<p$. Aber $4aa-p$ hat die Form $4n+3$, und $+p$ ist ein quadratischer Rest von $4aa-p$ (denn $p\equiv4aa\ (\mathrm{mod.}\,4aa-p)$). Ist also $4aa-p$ eine Primzahl, so ist $-p$ ein Nichtrest davon. Andernfalls muss $4aa-p$ einen Faktor der Form $4n+3$ haben. Da $+p$ auch von diesem ein Rest sein muss, ist $-p$ ein Nichtrest davon. \mbox{\textit{Q. E. D.}}

Für \textit{positiv} zu nehmende Primzahlen unterscheiden wir zwei Fälle. \textit{Erstens:} Sei $p$ eine Primzahl der Form $8n+5$. Sei $a$ eine beliebige positive Zahl $<\surd\!\frac12 p$. Dann ist $8n+5-2aa$ eine positive Zahl der Form $8n+5$ oder $8n+3$, je nachdem, ob $a$ \SourceLeaf{100}gerade oder ungerade ist. Sie ist daher durch eine Primzahl der Form $8n+3$ oder $8n+5$ teilbar; denn ein Produkt beliebig vieler Zahlen der Formen $8n+1$ und $8n+7$ kann weder die Form $8n+3$ noch die Form $8n+5$ haben. Diese Primzahl sei $q$; dann gilt $8n+5\equiv2a^2\ (\mathrm{mod.}\,q)$. Aber $2$ ist ein Nichtrest von $q$ (Artikel $112$), also auch $2a^2$\HistoricalNote{*)}{Artikel $98$. Denn $a^2$ ist ein Rest von $q$, der nicht durch $q$ teilbar ist; sonst wäre auch die Primzahl $p$ durch $q$ teilbar. \mbox{\textit{Q. E. A.}}} und $8n+5$. \mbox{\textit{Q. E. D.}}

\Article{126}
Dass jede positiv genommene Primzahl der Form $8n+1$ ein Nichtrest einer kleineren Primzahl ist, lässt sich aber nicht durch so naheliegende Kunstgriffe beweisen. Diese Wahrheit ist von größter Bedeutung; deshalb können wir einen strengen Beweis nicht übergehen, auch wenn er etwas länger ist. Wir stellen folgendes Lemma voran.

\Needspace{10\baselineskip}
\textsc{Lemma.} \textit{Es seien zwei Zahlenfolgen}

\[
A,B,C\ \textit{usw.}\ ....\ (\mathrm{I}),\quad A^{\prime},B^{\prime},C^{\prime}\ \textit{usw.}\ ....\ (\mathrm{II})
\]

(ob sie dieselbe Anzahl von Gliedern haben, spielt keine Rolle) \textit{so beschaffen, dass für jede Primzahl oder Primzahlpotenz $p$, die ein Glied oder mehrere Glieder der zweiten Folge teilt, mindestens ebenso viele Glieder der ersten Folge durch $p$ teilbar sind wie Glieder der zweiten. Dann behaupte ich, dass das Produkt aller Zahlen} (I) \textit{durch das Produkt aller Zahlen} (II) \textit{teilbar ist.}

\textit{Beispiel.} (I) bestehe aus den Zahlen $12,18,45$, (II) aus $3,4,5,6,9$. Dann sind durch $2,4,3,9,5$ in (I) jeweils $2,1,3,2,1$ Glieder, in (II) jeweils $2,1,3,1,1$ Glieder teilbar. Das Produkt aller Glieder von (I), $=9720$, ist durch das Produkt aller Glieder von (II), $3240$, teilbar.

\Needspace{16\baselineskip}
\textit{Beweis.} Das Produkt aller Glieder von (I) sei $=Q$, das aller Glieder von (II) $=Q'$. Jede Primzahl, die $Q'$ teilt, teilt offensichtlich auch $Q$. Wir zeigen nun, dass jeder Primfaktor von $Q'$ in $Q$ einen mindestens ebenso großen Exponenten hat wie in $Q'$. Sei $p$ ein solcher Teiler. In der Folge (I) seien $a$ Glieder durch $p$, $b$ Glieder durch $p^2$, $c$ Glieder durch $p^3$ teilbar usw.; die Buchstaben $a',b',c'$ usw. bedeuten Entsprechendes für (II). Dann hat $p$ in $Q$ offensichtlich den Exponenten \SourceLeaf{101}$a+b+c+$ usw., in $Q'$ dagegen $a'+b'+c'+$ usw. Nach Voraussetzung ist aber $a'$ nicht größer als $a$, $b'$ nicht größer als $b$ usw. Daher ist $a'+b'+c'$ usw. sicher nicht $>a+b+c$ usw. --- Da also keine Primzahl in $Q'$ einen größeren Exponenten haben kann als in $Q$, ist $Q$ durch $Q'$ teilbar (Artikel $17$). \mbox{\textit{Q. E. D.}}

\Needspace{10\baselineskip}
\Article{127}
\textsc{Lemma.} \textit{In der Folge $1,2,3,4,....n$ können nicht mehr Glieder durch eine beliebige Zahl $h$ teilbar sein als in der Folge $a,a+1,a+2.....a+n-1$, die ebenso viele Glieder hat.}

Ist $n$ ein Vielfaches von $h$, so sind in beiden Folgen offensichtlich $\frac nh$ Glieder durch $h$ teilbar. Andernfalls setze man $n=eh+f$, wobei $f$ $<h$ ist. Dann sind in der ersten Folge $e$ Glieder durch $h$ teilbar, in der zweiten entweder ebenso viele oder $e+1$.

Als Folgerung ergibt sich der aus der Theorie der figurierten Zahlen bekannte, aber, wenn wir uns nicht irren, bisher von niemandem direkt bewiesene Satz, dass

\[
\frac{a.a+1.a+2....a+n-1}{1\ .\ 2\ .\ 3\ ....\ n}
\]

stets eine ganze Zahl ist.

Schließlich hätte man das Lemma allgemeiner so aussprechen können:

In der Folge $a,a+1,a+2,....a+n-1$ gibt es mindestens ebenso viele Glieder, die nach dem Modul $h$ einer beliebigen gegebenen Zahl $r$ kongruent sind, wie in der Folge $1,2,3....n$ Glieder durch $h$ teilbar sind.

\Needspace{12\baselineskip}
\Article{128}
\textsc{Satz.} \textit{Sei $a$ eine beliebige Zahl der Form $8n+1$, $p$ eine beliebige zu $a$ teilerfremde Zahl, von der $+a$ ein Rest ist, und schließlich $m$ beliebig. Dann behaupte ich, dass in der Folge}

\[
\begin{gathered}a,\ \tfrac12(a-1),\ 2(a-4),\ \tfrac12(a-9),\ 2(a-16),\\ ....\ 2(a-m^2),\quad\textit{oder}\quad\tfrac12(a-m^2)\end{gathered}
\]

\textit{je nachdem, ob $m$ gerade oder ungerade ist, mindestens ebenso viele Glieder durch $p$ teilbar sind wie in der Folge}

\[
1,\ 2,\ 3,\ ....\ 2m+1
\]

Die erste Folge bezeichnen wir mit (I), die zweite mit (II).

\SourceLeaf{102}\textit{Beweis.} I. Ist $p=2$, so sind in (I) alle Glieder außer dem ersten, \textit{das heißt} $m$ Glieder, teilbar; ebenso viele sind es in (II).

II. Sei $p$ eine ungerade Zahl oder das Doppelte oder Vierfache einer ungeraden Zahl, und es gelte $a\equiv rr\ (\mathrm{mod.}\,p)$. Dann gibt es in der Folge $-m,-(m-1),-(m-2),.....+m$ (die ebenso viele Glieder wie $(\mathrm{II})$ hat und mit (III) bezeichnet wird) mindestens ebenso viele Glieder, die nach dem Modul $p$ der Zahl $r$ kongruent sind, wie in (II) Glieder durch $p$ teilbar sind (voriger Artikel). Unter diesen können aber keine zwei vorkommen, die sich nur im Vorzeichen, nicht im Betrag unterscheiden\HistoricalNote{*)}{Wäre nämlich $r\equiv-f\equiv+f\,(\operatorname{mod.}p)$, so wäre $2f$ durch $p$ teilbar und damit auch $2a$ (wegen $ff\equiv a\,(\operatorname{mod.}p)$). Das ist nur möglich, wenn $p=2$, da $a$ nach Voraussetzung zu $p$ teilerfremd ist. Diesen Fall haben wir aber schon gesondert behandelt.}. Jedem davon entspricht schließlich ein durch $p$ teilbares Glied von (I). Ist nämlich $\pm b$ ein Glied von (III), das $r$ nach $p$ kongruent ist, so ist $a-bb$ durch $p$ teilbar. Ist also $b$ gerade, so ist das Glied $2(a-bb)$ von (I) durch $p$ teilbar. Ist dagegen $b$ ungerade, so ist das Glied $\frac12(a-bb)$ durch $p$ teilbar. Denn $\frac{a-bb}{p}$ ist offensichtlich eine \textit{gerade} ganze Zahl, da $a-bb$ durch $8$ teilbar ist, $p$ aber durch keine höhere Zweierpotenz als $4$ teilbar ist (denn $a$ hat nach Voraussetzung die Form $8n+1$, und $bb$ als Quadrat einer ungeraden Zahl dieselbe Form; die Differenz hat also die Form $8n$). Somit gibt es in (I) ebenso viele durch $p$ teilbare Glieder wie in (III) Glieder, die $r$ nach $p$ kongruent sind, \textit{das heißt} mindestens ebenso viele, wie in (II) durch $p$ teilbar sind. \mbox{\textit{Q. E. D.}}

III. Sei $p$ von der Form $8n$, und es gelte $a\equiv rr\ (\mathrm{mod.}\,2p)$. Denn $a$ ist, da es nach Voraussetzung ein Rest von $p$ ist, offensichtlich auch ein Rest von $2p$. Dann gibt es in (III) mindestens ebenso viele Glieder, die $r$ nach $p$ kongruent sind, wie in (II) durch $p$ teilbar sind; jene haben alle verschiedene Beträge. Jedem davon entspricht aber ein durch $p$ teilbares Glied von (I). Denn wenn $+b$ oder $-b\equiv r\ (\mathrm{mod.}\,p)$, so gilt $bb\equiv rr\ (\mathrm{mod.}\,2p)$\HistoricalNote{\(\dagger\))}{Denn $bb-rr=(b-r)(b+r)$ ist ein Produkt aus zwei Faktoren, von denen einer durch $p$ teilbar ist (Voraussetzung), der andere durch $2$ (da sowohl $b$ als auch $r$ ungerade sind). Also ist $bb-rr$ durch $2p$ teilbar.}, und damit ist das Glied $\frac12(a-bb)$ durch $p$ teilbar. Daher gibt es in (I) mindestens ebenso viele durch $p$ teilbare Glieder wie in (II). \mbox{\textit{Q. E. D.}}

\Needspace{10\baselineskip}
\Article{129}
\textsc{Satz.} \textit{Ist $a$ eine Primzahl der Form $8n+1$, so gibt es unterhalb von $2\surd a+1$ notwendig eine Primzahl, von der $a$ ein Nichtrest ist.}

\SourceLeaf{103}\textit{Beweis.} Nehmen wir, wenn möglich, an, dass $a$ ein Rest aller Primzahlen unterhalb von $2\surd a+1$ ist. Dann ist $a$ auch ein Rest aller zusammengesetzten Zahlen unterhalb von $2\surd a+1$ (siehe die Regeln zur Entscheidung, ob eine Zahl ein Rest einer zusammengesetzten Zahl ist, Artikel $105$). Die ganze Zahl unmittelbar unterhalb von $\surd a$ sei $=m$. Dann gibt es in der Folge

\[
\begin{gathered}(\mathrm I)......a,\tfrac12(a-1),2(a-4),\tfrac12(a-9)....2(a-mm)\\ \text{oder}\quad\tfrac12(a-mm)\end{gathered}
\]

mindestens ebenso viele Glieder, die durch eine beliebige Zahl unterhalb von $2\surd a+1$ teilbar sind, wie in der Folge

\[
(\mathrm{II}).....1,2,3,4....2m+1\quad\text{(voriger Artikel)}
\]

Daraus folgt, dass das Produkt aller Glieder von (I) durch das Produkt aller Glieder von (II) teilbar ist (Artikel $126$). Das erste Produkt ist entweder \mbox{$=a(a-1)(a-4)....(a-mm)$} oder die Hälfte davon, je nachdem, ob $m$ gerade oder ungerade ist. Somit ist das Produkt $a(a-1)(a-4)....(a-mm)$ sicher durch das Produkt aller Glieder von $(\mathrm{II})$ teilbar. Da alle diese Glieder zu $a$ teilerfremd sind, gilt dies auch, wenn man den Faktor $a$ weglässt. Das Produkt aller Glieder von $(\mathrm{II})$ lässt sich aber auch schreiben als

\[
(m+1).((m+1)^2-1).((m+1)^2-4)....((m+1)^2-m^2)
\]

\Needspace{8\baselineskip}
Also ist

\[
\frac1{m+1}.\frac{a-1}{(m+1)^2-1}.\frac{a-4}{(m+1)^2-4}....\frac{a-m^2}{(m+1)^2-m^2}
\]

eine ganze Zahl, obwohl es ein Produkt von Brüchen kleiner als eins ist. Denn $\surd a$ muss irrational sein. Daher gilt $m+1>\surd a$ und damit $(m+1)^2>a.$ Unsere Annahme ist also unmöglich. \mbox{\textit{Q. E. D.}}

Da $a$ sicher $>9$ ist, gilt $2\surd a+1<a$. Somit gibt es eine Primzahl $<a$, von der $a$ ein Nichtrest ist.
